\documentclass[a4paper]{article}
\usepackage[affil-it]{authblk}
\usepackage{graphicx}
\usepackage{amsfonts}
\usepackage{amsmath}
\usepackage[backend=bibtex,style=numeric]{biblatex}

\usepackage{geometry}
\geometry{margin=1.5cm, vmargin={0pt,1cm}}
\setlength{\topmargin}{-1cm}
\setlength{\paperheight}{29.7cm}
\setlength{\textheight}{25.3cm}

\addbibresource{citation.bib}

\begin{document}
% =================================================
\title{Numerical Analysis homework 1}

\author{Zirui Chen 22435036
  \thanks{Electronic address: \texttt{18107732576@163.com}}}
\affil{(24), Zhejiang University }


\date{Due time: \today}

\maketitle

%\begin{abstract}
 %   The abstract is not necessary for the theoretical homework, 
   % but for the programming project, 
    %you are encouraged to write one.      
%\end{abstract}





% ============================================
\section*{I}

\subsection*{I-a}
prove:

we know the lagrange interpolation of f is 
$p_{1}(f;x)=f(x_{0})\frac{x-x_{1}}{x_{0}-x_{1}}+f({x_{1}})\frac{x-x_{0}}{x_{1}-x_{0}}=\frac{3}{2}-\frac{x}{2}$

then we can use the number of 
$x_{0},x_{1}$,$f(x)-p_{1}(f;x)=\frac{1}{2x}(x-1)(x-2) \Rightarrow f^{''}(\xi(x))=\frac{1}{x} \Rightarrow \xi(x)=\sqrt[3]{2x} $ 
\begin{flushright} % 右对齐
    $\square$ 
\end{flushright} 
\subsection*{I-b}
Prove:

Follows from above ,we have:
$\xi(x)=\sqrt[3]{2x} ,\max_{x \in [1,2]} \xi(x)=\sqrt[3]{4} ,\min_{x \in [1,2]} \xi(x)=\sqrt[3]{2} $
$,\max_{x \in [1,2]} f^{''}(\xi(x))=1 $
\begin{flushright} 
    $\square$ 
\end{flushright} 


\section*{II}
Prove:

Use lagrange interpolation 
$p_{n}{f;x}=\sum_{k = 0}^{n}f_{k}(x)(\frac{\pi _{n+1}(x)}{(x-x_{k})\pi^{'}_{n+1}(x_{k})})^{2}$

it is obvious that $p_{n}(f;x) \in \mathbb{P}_{2n}^{+} $
\begin{flushright} 
    $\square$ 
\end{flushright} 


\section*{III}
\subsection*{III-a}
Prove:

To prove the induction ,we begin with the base case $n=0$,$f[t]=f(t)=e^{t}$

Assume that the induction hypothesis is true for $n=k-1$,We need to show that the statement holds for $n=k$.

We have:

\begin{align}
    f[t,t+1,\cdots,t+n] & =\frac{1}{t+n-t}(f[t+n,\cdots,t+1]-f[t+n-1,\cdots,t])\\
                        & =\frac{1}{n}(\frac{(e-1)^{n-1}}{(n-1)!}e^{t+1}-\frac{(e-1)^{n-1}}{(n-1)!}e^{t})\\
                        & =\frac{(e-1)^{n-1}}{n!}e^{t}
\end{align}


\begin{flushright} 
    $\square$ 
\end{flushright} 

\subsection*{III-b}
Prove:

we know $f^{(n)}(x)=e^{x}$and $f[t,t+1,\cdots,t+n]=\frac{(e-1)^{n}}{n!}e^{t}$ .

$ \Rightarrow f[0,1,\cdots,n]=\frac{(e-1)^{n}}{n!} \Rightarrow \xi=nln(e-1) $ ,
because $e-1>\sqrt{e} \Rightarrow \xi>\frac{n}{2}$ .
\begin{flushright} 
    $\square$ 
\end{flushright} 


\section*{IV}

\subsection*{IV-a}
Newton fomular:

\begin{align}
    p_{3}(f;x) & =f[0]+f[0,1](x-0)+f[0,1,3]+x(x-1)+f[0,1,3,4]x(x-1)(x-3)\\
                 &  =5-2x+x(x-1)+\frac{1}{4}x(x-1)(x-3).
\end{align}

\subsection*{IV-b}
Find minimun :

we just need to find the minimum of $p_{3}(f;x)$,then $x_{min}=\sqrt{5}$.
\begin{flushright} 
    $\square$ 
\end{flushright} 


\section*{V}

\subsection*{V-a}
Compute f[0,1,1,1,2,2] :


\begin{table}[h]
    \centering
    \begin{tabular}{c|cccccc}
        
        0 &0 & & & & & \\
        
        1 &1 &1 & & & & \\
        
        1 &1 &7 &6 & & & \\
        
        1 &1 &7 &21 &15 & & \\
        
        2 &128 &127 &120 &99 &42 & \\
        
        2 &128 &448 &321 &201 &102 &30 \\
    \end{tabular} 
\end{table}
So ,$f[0,1,1,1,2,2] = 30$ .


\subsection*{V-b}
Determine $\xi$ :

Newton interpolation of $x^{7}$:$p(f;x)=0+x+6x(x-1)+15x(x-1)(x-1)^{2}+42x(x-1)^{3}+30x(x-1)^{3}(x-2)+\frac{f^{(6)}(\xi)}{6!}$.

It follows from Corollary2.22 that $f[x_{0}.x_{1}.\cdots.x_{n}]=\frac{1}{n!}f^{(n)}(\xi) \Rightarrow \xi= 1.195$ .

\section*{VI}

\subsection*{VI-a}
Estimate :

First we give the Hermite interpolation formula :
$$p_{4}(f;x)=1+x-2x(x-1)+\frac{2}{3}x(x-1)^{2}-\frac{5}{36}x(x-1)^{2}(x-3)$$,
Then we get $f(2)\approx p_{4}(f;2)=\frac{11}{18}.$

\subsection*{VI-b}
Estimate :

Follows from the Theorem 2.37 we have ,
\begin{align}
    &f(x)-p_{4}(f;x) =\frac{f^{5}(\xi)}{5!}(x-1)(x-2)^{2}(x-3) ,\\
    \Rightarrow &f(2)-p_{4}(f;2)  = \frac{f^{5}(\xi)}{5!} \leq \frac{M}{60}.
\end{align}

\section*{VII}
Prove :

Assume $x_{0}=x $.$x_{1}=x+h$ ,the induction basis holds for $n=1$ ,because
$$\bigtriangleup f(X)=f(x+h)-f(x)=f[x_{1}]-f[x_{0}]=hf[x_{0},x_{1}]$$.
Suppose the induction hypothesis holds for $n-1$, we need to show that it holds for $n$.
\begin{align}
    \bigtriangleup ^{n} f(x)&= \bigtriangleup^{n-1} \bigtriangleup f(x)=\bigtriangleup ^{n-1}(f(x+h)-f(x)) \\
    &=\bigtriangleup ^{n-1}f(x+h)-\bigtriangleup ^{n-1}f(x)=(n-1)!h^{n-1}(f[x_{1},\cdots,x_{n}]-f[x_{0},\cdots,x_{n-1}])\\
    &=(n-1)!\cdot h^{n-1}(x_{n}-x_{0})f[x_{0},\cdots,x_{n}]= n!\cdot h^{n}f[x_{0},\cdots,x_{n}].
\end{align}
The similar result holds for $\bigtriangledown ^{n} f(x)=n!\cdot h^{n}f[x_{0},\cdots,x_{n-k}]$.
\begin{flushright} 
    $\square$ 
\end{flushright} 

\section*{VIII}
Prove :

Notice that 
\begin{align}
    \frac{\partial}{ \partial x_{0}}f[x_{0},\cdots,x_{n}] &=  
    \lim_{\bigtriangledown x \to 0} \frac{1}{\bigtriangledown x}(f[x_{0}+\bigtriangledown x,\cdots,x_{n}]-f[x_{0},\cdots,x_{n}])\\
    &= \lim_{\bigtriangledown x \to 0}f[x_{0}+\bigtriangledown x,x_{0},\cdots,x_{n}]\\    
    &=f[x_{0},x_{0},\cdots,x_{n}].
\end{align}
It is not sure about other variables ,because we dont know the differential of the variable.
\begin{flushright} 
    $\square$ 
\end{flushright} 


\section*{IX}
Determine :

First we denote $p(x)$ as the polynomial which the problem asked ,
consider $x=(\frac{b-a}{2})y+\frac{a+b}{2}$ ,which $y\in [-1,1]$.Than the polynomial $p(y)$ has the leading coefficient
$a_{0}(\frac{b-a}{2})^{n}$ ,a chebyshev polynomial of degree $n$ in $[-1,1]$ has maximun number 1.
Then we follow from the Theorem 2.47 to get the minimum maximun value of $p(y)$ is $a_{0}\frac{(b-a)^{n}}{2^{2n-1}}$.
It is also the minimum maximun value of $p(x)$.


\section*{X}
Prove :

Transform the chebyshev polynomial's domain from $[-1,1]$ to $\mathbb{R} $ ,we have:
$$T_{0}(x)=1\ ,T_{1}(x)=x\ ,T_{2}(x)=2x^{2}-1\ ,T_{3}(x)=4x^{3}-3x\ ,\cdots\ ,T_{n+1}(x)=2xT_{n}(x)-T_{n-1}(x)$$,
if $\exists p \in \mathbb{P} _{n}^{a} $ ,such that $\|p\|_{\infty }<\|\widehat{p}_{n}(x)\|_{\infty } $,
Consider $Q(x)=\widehat{p}_{n}(x)-p(x) $ ,then we know that $Q(x)$ is a polynomial which has n zero points in $[-1,1]$ ,
and it is obvious that $Q(a)=0$.But it constrat to Q(x) is a n degree polynomial in $\mathbb{R} $.
\begin{flushright} 
    $\square$ 
\end{flushright} 


\section*{XI}
Prove :

Use the porperty of Berstain polynomials,
\begin{align}
    \frac{n-k}{n}b_{n,k}(t)+\frac{k+1}{n}b_{n,k+1}(t) 
    &= \frac{n-k}{n}\binom{n}{k}t^{k}(1-t)^{n-k}+\frac{k+1}{n}\binom{n}{k+1}t^{k+1}(1-t)^{n-k-1}\\
    &=t^{k}(1-t)^{n-k-1}(\frac{n-k}{n}\frac{n!}{k!(n-k)!}(1-t)+\frac{k+1}{n}\frac{n!}{(k+1)!(n-k-1)!}t )\\
    &=t^{k}(1-t)^{n-k-1}\binom{n-1}{k}\\
    &=b_{n-1,k}(t).
\end{align}
\begin{flushright} 
    $\square$ 
\end{flushright} 

\section*{XII}
Prove :

For $\forall k  = 0,1,2,\cdots,n$, we have ,$\int _{0}^{1}b_{n,k}(t)dt=\frac{1}{n+1}$.

\begin{align}
    \int _{0}^{1}\frac{n!}{k!(n-k)!}t^{k}(1-t)^{n-k}dt&=\frac{n!}{k!(n-k)!}\cdot\frac{n-k}{k+1}\int _{0}^{1} t^{k+1}(1-t)^{n-k-1}dt\\
    &=\frac{n!}{k!(n-k)!}\cdot\frac{(n-k)(n-k-1)}{(k+1)(k+2)}\int _{0}^{1}t^{k+2}(1-t)^{n-k-2}dt\\
    &=\cdots=\frac{n!}{k!(n-k)!}\cdot\frac{(n-k)!k!}{n!}\int _{0}^{1} t^{n}dt\\
    &=\frac{1}{n+1}
\end{align}
\begin{flushright} 
    $\square$ 
\end{flushright} 


% ===============================================
\section*{ \center{\normalsize {Acknowledgement}} }
Give your acknowledgements here(if any).


\printbibliography

If you are not familiar with \texttt{biblatex}, 
\cite{Zhang2020BooleanAO}
it is acceptable to put a table here for your references.
\end{document}